% ---------------------------------------------------------------------------
% Préambule — rapport de clôture JUNON
% Compilation : pdflatex + biber (voir Makefile)
% ---------------------------------------------------------------------------

\usepackage[T1]{fontenc}
\usepackage[utf8]{inputenc}
\usepackage{libertinus}          % corps à empattements, usage académique
\usepackage[french]{babel}
\usepackage{microtype}

% --- Mise en page ---------------------------------------------------------
\usepackage[a4paper,top=2.6cm,bottom=2.6cm,left=2.6cm,right=2.6cm,headsep=1.1em]{geometry}
\usepackage{setspace}
\setstretch{1.08}

% --- Couleurs : sobriété, un seul accent pour les titres et les liens ------
\usepackage{xcolor}
\definecolor{encre}{HTML}{1A1A1A}
\definecolor{accent}{HTML}{243B53}
\definecolor{gris}{HTML}{555555}
\definecolor{grisfin}{HTML}{BFBFBF}

% --- Titres ---------------------------------------------------------------
\setkomafont{disposition}{\normalfont\bfseries\color{accent}}
\setkomafont{chapter}{\normalfont\Large\bfseries\color{accent}}
\setkomafont{section}{\normalfont\large\bfseries\color{accent}}
\setkomafont{subsection}{\normalfont\normalsize\bfseries\color{encre}}
\setkomafont{caption}{\small}
\setkomafont{captionlabel}{\small\bfseries}

\RedeclareSectionCommand[beforeskip=-1.3\baselineskip, afterskip=0.7\baselineskip]{chapter}

% Titre de chapitre : numéro détaché, filet de séparation sous le titre.
\renewcommand*{\chapterformat}{%
  \textcolor{grisfin}{\Large\bfseries\thechapter}\enskip
}
\renewcommand*{\chapterlinesformat}[3]{%
  \parbox{\textwidth}{%
    \raggedright #2#3%
    \par\vspace{0.35em}%
    \noindent\textcolor{grisfin}{\rule{\textwidth}{0.6pt}}%
  }%
}
\RedeclareSectionCommand[beforeskip=-1.8ex plus -0.4ex minus -0.2ex,
                         afterskip=0.9ex plus 0.2ex]{section}
\RedeclareSectionCommand[beforeskip=-1.5ex plus -0.3ex minus -0.2ex,
                         afterskip=0.7ex plus 0.2ex]{subsection}

% --- En-têtes -------------------------------------------------------------
\usepackage[automark,headsepline]{scrlayer-scrpage}
\pagestyle{scrheadings}
\clearpairofpagestyles
\ihead{\small\itshape\color{gris}\headmark}
\ohead{\small\color{gris}\thepage}
\setkomafont{pageheadfoot}{\normalfont\small}

% --- Tableaux -------------------------------------------------------------
\usepackage{booktabs}
\usepackage{tabularx}
\usepackage{longtable}
\usepackage{array}
\renewcommand{\arraystretch}{1.2}
\setlength{\heavyrulewidth}{0.9pt}
\setlength{\lightrulewidth}{0.4pt}
\setlength{\aboverulesep}{0.5ex}
\setlength{\belowrulesep}{0.6ex}

% --- Figures --------------------------------------------------------------
\usepackage{graphicx}
\usepackage{float}
\usepackage{tikz}
\usetikzlibrary{arrows.meta,positioning,fit,backgrounds,calc}

% ---------------------------------------------------------------------------
% Mise en relief : pas de fond coloré ni de cadre. Un filet vertical fin et un
% retrait suffisent à distinguer une remarque du corps du texte.
% ---------------------------------------------------------------------------
\usepackage[most]{tcolorbox}

\newtcolorbox{remarque}[1][]{
  enhanced, breakable,
  colback=white, colframe=grisfin,
  boxrule=0pt, leftrule=1.1pt,
  sharp corners, arc=0pt,
  left=10pt, right=2pt, top=3pt, bottom=3pt,
  fontupper=\small,
  coltitle=accent, fonttitle=\small\bfseries,
  attach title to upper={\par\nobreak\vspace{2pt}},
  #1
}

\newtcolorbox{mesure}[1][]{
  enhanced, breakable,
  colback=white, colframe=accent,
  boxrule=0pt, leftrule=1.6pt,
  sharp corners, arc=0pt,
  left=10pt, right=2pt, top=3pt, bottom=3pt,
  fontupper=\small,
  coltitle=accent, fonttitle=\small\scshape,
  attach title to upper={\par\nobreak\vspace{2pt}},
  #1
}

% --- Question d'ouverture de chapitre -------------------------------------
\newcommand{\questionchapitre}[1]{%
  \begingroup
  \small\itshape\color{gris}
  \begin{quote}
  #1
  \end{quote}
  \endgroup
  \vspace{0.2em}
}

% --- Bilan de fin de chapitre ---------------------------------------------
\newenvironment{acquisouvert}{%
  \par\vspace{0.9em}
  \noindent\textcolor{grisfin}{\rule{\textwidth}{0.5pt}}
  \par\vspace{0.4em}
  \small
}{%
  \par\vspace{0.9em}
}

% --- Listes ---------------------------------------------------------------
\usepackage{enumitem}
\setlist[itemize]{leftmargin=1.3em,itemsep=0.2em,topsep=0.35em}
\setlist[enumerate]{leftmargin=1.5em,itemsep=0.2em,topsep=0.35em}

% --- Unités ---------------------------------------------------------------
\usepackage{siunitx}
\sisetup{
  output-decimal-marker = {,},
  group-separator       = {\,},
  group-minimum-digits  = 4,
  range-phrase          = {\ à\ },
  range-units           = single,
}
\DeclareSIUnit{\an}{an}
\DeclareSIUnit{\euro}{\text{\texteuro}}

% --- Bibliographie --------------------------------------------------------
\usepackage[backend=biber,style=numeric,sorting=nyt,maxbibnames=8,giveninits=true,
            doi=false,isbn=false,url=true]{biblatex}
\addbibresource{bibliographie.bib}

% --- Liens ----------------------------------------------------------------
\usepackage{hyperref}
\hypersetup{
  colorlinks=true,
  linkcolor=accent,
  citecolor=accent,
  urlcolor=accent,
  pdftitle={Rapport de clôture, programme JUNON},
  pdfauthor={Nicolas Ringuet},
}
\usepackage{url}
\urlstyle{same}

% --- Divers ---------------------------------------------------------------
\usepackage{csquotes}
\usepackage{caption}
\captionsetup{font=small,labelfont=bf}

% --- Raccourcis -----------------------------------------------------------
\newcommand{\code}[1]{\texttt{\small #1}}
\newcommand{\hubeau}{\code{hubeau\_data\_integration}}
\newcommand{\junonrepo}{\code{time-serie-explo}}
